\begin{enumerate}
    \item La formule de l'ion fer(III) s'écrit $\ce{Fe^{3+}}$. La neutralité
          électrique impose trois fois plus d'ions chlorure que d'ions fer(III).

          La formule statistique du chlorure de fer(III) s'écrit donc
          \ce{FeCl3}.
    \item La quantité de matière que l'on doit dissoudre se calcule à partir de
          la formule suivante~:
          \[
              c=\frac{n_{\text{soluté}}}{V_{\text{solution}}}
          \]

          Soit~:
          $n_{\ce{FeCl3}}=c\times V=0,250\times0,100=\qty{2.50e-2}{\mol}$.
    
          La masse molaire du chlorure de fer(III) vaut~:

          \[
              M_{\ce{FeCl3}}=55,8+3\times35,5=\qty{162}{\g\per\mol}
          \]

          Il faut donc peser une masse de chlorure de fer(III) égale à~:

          \[
              m=n\times M=2,50\times 10^{-2}\times 162\approx\qty{4.05}{\g}.
          \]

    \item L'équation de la réaction de dissolution du chlorure de fer(III)
          s'écrit~:
          \[
              \ce{FeCl3(s) -> Fe^{3+}(aq) + 3Cl-(aq)}.
          \]
          Si, dans un volume $V$, on a dissous 1 mole \ce{FeCl3(s)}, on obtient
          en solution 1 mole \ce{Fe^{3+}(aq)} et 3 moles \ce{Cl-(aq)}. Donc~:
          \begin{align*}
              &\left[\ce{Fe^{3+}(aq)}\right]=c=\qty{0.250}{\mol\per\L} \\
              &\left[\ce{Cl-(aq)}\right]=3c=\qty{0.750}{\mol\per\L}.
          \end{align*}
\end{enumerate}



